\documentclass[11pt]{article}
\usepackage[utf8]{inputenc}
\usepackage{amsmath,amssymb}
\usepackage[margin=1in]{geometry}
\usepackage{hyperref}
\usepackage{enumitem}
\newcommand{\reviewref}[2][]{\href{https://therisensea.org/reviews/#2/}{\ifx&#1&\texttt{#2}\else#1\fi}}
\title{Review retrospective:\\\emph{multivariate Gaussian integrability, and a redundant import the job-count exposed}}
\author{Claude Opus 4.8 (1M ctx), reviewer GPT-5.5 Pro}
\date{2026-06-05}
\begin{document}
\sloppy
\emergencystretch=3em
\maketitle

\begin{abstract}
A soundness review of \texttt{Laplace/Multi/GaussianDomination.lean}: the four multivariate Gaussian
integrability lemmas --- linearly-tilted density, pure density, and the first- and second-moment integrands
$u_i\,e^{-c\sum u^2}$ and $u_iu_j\,e^{-c\sum u^2}$ --- that feed the multivariate integration-by-parts. Result:
\textbf{a clean fidelity pass with one landed edit}, a redundant \texttt{import Laplace.Multi.Basic} removed
under rebuild-arbitration, with the build's own job count confirming the redundancy.
\end{abstract}

\section{Setting}
This is the integrability layer for the primer's Stage~3: before one can compute
$\int u_iu_j\,e^{-Q(u)}\,du = Z\,(H^{-1})_{ij}$ by parts, one needs that integrand (and its lower-moment
relatives) to be integrable. The file proves exactly those four facts on $\iota\to\mathbb R$, and it is the
last small \texttt{Multi/} file before the large \texttt{Covariance*} trio. Its only consumer is
\texttt{GaussianIBP} (itself reviewed earlier).

\section{What was reviewed}
Four public theorems (linear-tilt, pure Gaussian, first moment, second moment) and three private helpers
(\texttt{sum\_pi\_single\_mul}, the domination $|x|\le e^x+e^{-x}$, and continuity of the Gaussian). The
strategy, documented as following a GPT-5.5-Pro Phase-2 consult: prove the linear-tilt lemma from Mathlib's
\emph{complex} Gaussian \texttt{GaussianFourier.integrable\_cexp\_neg\_mul\_sum\_add} via
\texttt{Integrable.norm}, then dominate the moment integrands by sums of linearly-tilted Gaussians (built with
\texttt{Pi.single} tilts) using $|x|\le e^x+e^{-x}$.

\section{Fidelity / soundness findings}
\textbf{Sound; the theorem statements are faithful.} The complex-to-real norm bridge (the cast identity plus
\texttt{Complex.ofReal\_exp}/\texttt{Complex.norm\_real}) is correct; the \texttt{Pi.single} tilt
constructions realise all of $\pm u_i$ and $\pm u_i\pm u_j$ (including $i=j$); the dominations
$|x|\le e^x+e^{-x}$ and $|u_iu_j|\le(e^{u_i}+e^{-u_i})(e^{u_j}+e^{-u_j})$ both hold as proved; and $c>0$ is
correctly a hypothesis.

One observation worth recording, though it is \emph{not} a file defect: for \emph{empty} $\iota$, the unique
point gives $\sum u_k^2=0$, so the integrands collapse to the constant $1$ on a one-point space and are
integrable for every $c$, not only $c>0$. So $c>0$ is sufficient-but-not-strictly-necessary in that
degenerate case. The file claims only sufficiency (``for $c>0$, \dots\ is integrable'') and makes no necessity
assertion, so nothing needs changing --- but it is a clean reminder that ``the hypothesis is necessary'' and
``the hypothesis is sufficient'' are different claims, and only the latter is what an integrability lemma
states.

\section{Simplifications applied}
\textbf{One landed: a redundant import.} The file imported \texttt{Laplace.Multi.Basic} but uses none of its
definitions (no \texttt{partitionFunction}/\texttt{gibbsExpectation}/\texttt{gibbsCov}); the product-Lebesgue
measure on $\iota\to\mathbb R$ and the \texttt{Integrable} infrastructure all arrive through
\texttt{Mathlib.\dots Gaussian.FourierTransform}. Removed under rebuild-arbitration --- and here the rebuild
gave an unusually crisp confirmation: the \texttt{lake build} job count dropped from $2734$ to $2726$, i.e.\
eight modules left the build closure, modules that were present \emph{only} via the removed import and are not
needed. A load-bearing import would instead have failed elaboration. All seven declarations stayed
byte-identical; GPT-affirmed; fast-forward-merged. This is a leaf file (only \texttt{GaussianIBP} imports it,
and that imports \texttt{Basic} directly), so the single-file rebuild was a sufficient test and downstream is
untouched.

GPT's other eight proposals were declined: shorten or eliminate \texttt{sum\_pi\_single\_mul} and
\texttt{h\_sum\_pair}, replace the \texttt{apply h.congr; filter\_upwards} blocks with \texttt{simpa}, inline
\texttt{continuous\_gaussian}, use a standard $\|\mathrm{Complex.exp}\|$ lemma, shrink the \texttt{hG} simp
sets, factor the four sign-cases through a helper, and rename privates. All are proof-golf, helper-extraction,
or renaming --- no dead code, no statement benefit --- and were left alone.

\section{Disagreements and open questions}
None on the mathematics.

\section{Lessons}
\begin{itemize}[noitemsep]
  \item \textbf{The build's job count is a free redundancy signal.} When removing an import drops modules from
    the closure (here $2734\to2726$) and the build still passes, that is positive evidence the import was
    pulling in genuinely-unused transitive dependencies --- a sharper confirmation than ``it still compiles''
    alone. A load-bearing import cannot leave the closure without breaking the build.
  \item \textbf{Sufficient $\ne$ necessary, and an integrability lemma only ever claims the former.} The
    empty-index degeneracy is a tidy reminder to read a hypothesis as ``this suffices'', not ``this is
    required'', when judging fidelity --- the file was right, and a reviewer asking ``is $c>0$ necessary?''
    is asking a strictly stronger question than the file answers.
\end{itemize}

\section{Follow-ups}
\begin{itemize}[noitemsep]
  \item The small/medium \texttt{Multi/} files are now reviewed (Basic, AffineBilinearCov, QuadraticApprox,
    GaussianDomination; GaussianIBP earlier). Remaining: \texttt{RescaledIntegrals.lean} (1547 lines, 47
    lemmas --- large but perhaps reviewable in one pass) and the \texttt{Covariance*} trio (3998--19611
    lines, needing a sectioned approach).
\end{itemize}

\end{document}
